\scanpage{23}
% L1
\[
x_0\precy_0\mathrel{\overset{\text{déf}}{\Longleftrightarrow}}
\unsure{!}\,\exists z_1\in D_1\quad\text{t.q.}\quad x_0=\delta_0(z_1),y_0=\delta_1(z_1)
\]
% L1–4 left
% slanted
\margin{\deleted{\unsure{Bref, faut} introduire une \textquotedblleft fonction dimension\textquotedblright\ $\delta$, et exiger (pour $x_0<y_0$) que $\delta(x_0)<\delta(y_0)$, pour rendre la relation écrite \unsure{non-symétrique|anti-symétrique} !}}
% L2
La condition \uuline{est} $x_0<y_0$ dans le \textquotedblleft simplexe\textquotedblright\ $\unsure{\{x_0,y_0\}}$
\deleted{\unsure{C'est la}}
% L3–4
\deleted{A-t-on transitivité pour cette relation ? \unsure{Soit}}
\textcircled{$2^{\circ}$} La relation $\prec$ \deleted{\unsure{est dans|est bien}} sur $D_0$ est transitive.
\margin{?}
Ceci si on suppose satisfaite la condition suivante
% L5–7
\deleted{\textcircled{\unsure{2$'$}}}\inserted{\textcircled{2 bis}} L'application canonique
\[
D_2\xrightarrow{(\delta_2,\delta_0)}(D_1,\delta_1)\mathbin{\times}_{D_0}(D_1,\delta_0)
\]
est surjective,

% L8–13
Car si $x_0\prec y_0$ et $y_0\prec z_0$, i.e.
\[
x_0=\delta_0(z_1),y_0=\delta_1(z_1),\deleted{\unsure{x_0|y_0}}y_0=\delta_0(z'_1),z_0=\delta_1(z'_1),
\]
$\exists u_2\in D_2$ avec $z_1=\delta_2(u_2),z'_1=\delta_0(u_2)$, posons alors $z''_1=\delta_1(u_2)$, on a
\[
\begin{aligned}
\delta_0(z''_1)&=\delta_0\delta_1(u_2)=\delta_0\delta_2u_2\\
&=\delta_0(z_1)=x_0,\quad\delta_1(z''_1)=\delta_1\delta_1(u_2)=\delta_1\delta_0(u_2)=\delta_1(z'_1)=z_0,
\end{aligned}
\]
donc $x_0\prec z_0$, cqfd. \deleted{Je dis que la relation de \unsure{préor|préordre}}
% L14–15
Ainsi, posant
\[
x_0\unsure{\preceq|\leq|\preccurlyeq} y_0\mathrel{\overset{\text{déf}}{\Longleftrightarrow}}x_0=y_0\ \text{ou}\ x_0\prec y_0
\]
% L16–17
on trouve une relation de préordre dans $D_0$, \inserted{moyennant \textcircled{$2^{\circ}$}.}
\inserted{C'est une relation d'ordre, car $x_0\prec y_0$ et $y_0\prec x_0$ sont impossibles}
% L17–20 left
% slanted
\margin{sinon $x_0\prec x_0$ i.e. $\exists z_1$ avec $\delta_0(z_1)=\unsure{\delta_0|\delta_1}(z_1)=x_0$, impossible par \textcircled{1}. Cet ordre sur $D_0$ satisfait la condition \textcircled{b} Si $x_0\prec y_0$, $(x,y\in D_0)$ \deleted{\unsure{$\delta$|$\exists$}} $\exists z\in D_*$ tel que $x_0<z,y_0<z$}
% L18–21
Ceci \deleted{\unsure{dit}} posé, si $d\in D_*$, on voit que la relation d'ordre sur $I_d=\{x\in D_0\mid x<d\}$ est induite par la relation précédente. En effet, si $x,y\in D_0$, $x,y<d$ si on a
% L22–25
$x<y$ dans \deleted{\unsure{$\Delta_d,I$}} $I_d$, \unsure{alors} on a $x\prec y$, i.e. \inserted{l'inclusion} $I_d\hookrightarrow D_0$ est strictement croissante, d'où on déduit facilement \deleted{que} (l'ordre de $I_d$ étant total \inserted{et $D_0$ étant ordonné (\unsure{non pas seulement préord.!|non pas que préord.|voir paragraphe précédent})}) que l'ordre de $I_d$ est induit par celui de $D_0$.
% archive sheet number: 22
